% CARE-U User Documentation
% Compile: pdflatex user-documentation.tex (run twice for TOC)
\documentclass[11pt,a4paper]{report}

\usepackage[utf8]{inputenc}
\usepackage[T1]{fontenc}
\usepackage{lmodern}
\usepackage[margin=2.5cm]{geometry}
\usepackage{hyperref}
\usepackage{booktabs}
\usepackage{longtable}
\usepackage{enumitem}
\usepackage{xcolor}
\usepackage{tcolorbox}

\hypersetup{
    colorlinks=true,
    linkcolor=blue!60!black,
    urlcolor=blue!60!black,
    pdftitle={CARE-U User Guide},
}

\newtcolorbox{tipbox}{
    colback=blue!5,
    colframe=blue!40!black,
    title=\textbf{Tip},
    fonttitle=\bfseries,
}

\newtcolorbox{warningbox}{
    colback=orange!8,
    colframe=orange!60!black,
    title=\textbf{Important},
    fonttitle=\bfseries,
}

\title{\textbf{CARE-U}\\[0.4em]
\large User Guide}
\author{Hospital Staff \& Administrators}
\date{\today}

\begin{document}

\maketitle
\tableofcontents
\newpage

%==============================================================================
\chapter{Getting Started}
%==============================================================================

\section{What is CARE-U?}

CARE-U is a cloud hospital management system. Each hospital gets its own private workspace with modules for patients, appointments, clinical records, billing, pharmacy, laboratory, HR, and reports.

\section{Registering Your Hospital}

\begin{enumerate}
    \item Go to the platform website and click \textbf{Register}.
    \item Fill in:
    \begin{itemize}[noitemsep]
        \item Hospital name
        \item Hospital code (subdomain) --- e.g.\ \texttt{my-clinic}
        \item Admin name, email, phone
        \item Admin username and password
    \end{itemize}
    \item Submit the form. Your workspace is created automatically.
    \item On the success page, note your \textbf{Hospital code} --- you need this to sign in.
\end{enumerate}

\begin{warningbox}
Each hospital is completely separate. Data from one hospital is never visible to another. If you register a second hospital, you must sign in with \textbf{that hospital's code}.
\end{warningbox}

\section{Signing In}

Open \textbf{Sign In} and enter:

\begin{table}[h]
\centering
\begin{tabular}{@{}lp{8cm}@{}}
\toprule
\textbf{Field} & \textbf{What to enter} \\
\midrule
Hospital code & Your subdomain from registration (e.g.\ \texttt{my-clinic}) \\
Username or Email & The admin username you chose, or your email \\
Password & Your password \\
\bottomrule
\end{tabular}
\end{table}

After login you are taken to your hospital dashboard at:
\begin{center}
\texttt{/h/your-hospital-code/}
\end{center}

\begin{tipbox}
If your username is \texttt{admin} (common default), always enter your \textbf{hospital code}. Many hospitals use \texttt{admin} --- without the code, the system cannot tell which hospital you mean.
\end{tipbox}

\section{Completing Onboarding}

New hospitals should complete the onboarding wizard:
\begin{enumerate}[noitemsep]
    \item Set departments
    \item Choose brand colors
    \item Seed default catalog data (lab tests, drugs, services)
\end{enumerate}

%==============================================================================
\chapter{Dashboard}
%==============================================================================

The dashboard shows role-specific summary cards:

\begin{itemize}[noitemsep]
    \item \textbf{Admin / Receptionist} --- patients today, appointments, pending bills, revenue
    \item \textbf{Doctor} --- today's appointments, pending prescriptions
    \item \textbf{Receptionist} --- queue count
    \item \textbf{Pharmacist} --- pending dispenses, expiring stock
    \item \textbf{Lab technician} --- pending samples and results
\end{itemize}

Admins and accountants also see revenue and appointment charts.

%==============================================================================
\chapter{Patients}
%==============================================================================

\section{Registering a Patient}

\textbf{Who:} Receptionist, Admin

\begin{enumerate}[noitemsep]
    \item Sidebar $\rightarrow$ \textbf{Register Patient} (or Patients $\rightarrow$ Add)
    \item Enter CNIC (format: \texttt{12345-1234567-1}), name, phone, and optional demographics
    \item Save --- an MR number is generated automatically (e.g.\ \texttt{CARE-U-2026-00001})
\end{enumerate}

\section{Patient Profile}

Open any patient to see:
\begin{itemize}[noitemsep]
    \item Demographics and emergency contacts
    \item \textbf{Billing tab} --- invoices, total billed, paid, and balance due
\end{itemize}

\section{Bulk Import (CSV)}

\textbf{Who:} Admin, Receptionist

Sidebar $\rightarrow$ \textbf{Import Data} $\rightarrow$ Patients

\begin{enumerate}[noitemsep]
    \item Download the sample CSV
    \item Fill required columns: \texttt{cnic}, \texttt{first\_name}, \texttt{last\_name}, \texttt{phone}
    \item Upload and review import results
\end{enumerate}

%==============================================================================
\chapter{Appointments}
%==============================================================================

\section{Booking an Appointment}

\textbf{Who:} Receptionist, Admin

\begin{enumerate}
    \item Go to \textbf{Appointments} $\rightarrow$ \textbf{Book}
    \item Select patient, appointment type, \textbf{date}, and \textbf{doctor}
    \item The \textbf{Available time slot} dropdown shows only free slots based on:
    \begin{itemize}[noitemsep]
        \item Doctor's on-duty status
        \item Weekly schedule
        \item Date exceptions (leave, etc.)
        \item Already booked times
    \end{itemize}
    \item Choose a slot and confirm
\end{enumerate}

\begin{warningbox}
If no slots appear, the doctor may be off duty or has no schedule set. Ask the doctor to configure availability under \textbf{My Availability}.
\end{warningbox}

\section{Queue Board}

\textbf{Who:} Receptionist

\textbf{Queue} shows today's waiting patients. Actions:
\begin{itemize}[noitemsep]
    \item \textbf{Call} --- announce token
    \item \textbf{Complete} --- mark consultation done
\end{itemize}

\section{Doctor: My Availability}

\textbf{Who:} Doctor

\begin{itemize}[noitemsep]
    \item Toggle \textbf{On Duty / Off Duty}
    \item Add weekly schedule slots (day, start/end time)
    \item Add date exceptions for leave or special hours
\end{itemize}

When off duty, the doctor is hidden from appointment booking.

%==============================================================================
\chapter{Clinical}
%==============================================================================

\textbf{Who:} Doctors, Nurses, Admin

\begin{itemize}[noitemsep]
    \item \textbf{Visits} --- record OPD/IPD encounters
    \item \textbf{Vitals} --- temperature, BP, pulse, etc.
    \item \textbf{Wards} --- bed management for in-patients
    \item \textbf{Prescriptions} --- issued during visits
\end{itemize}

%==============================================================================
\chapter{Billing}
%==============================================================================

\section{Creating Invoices}

\textbf{Who:} Admin, Receptionist, Accountant

\begin{enumerate}[noitemsep]
    \item Billing $\rightarrow$ Invoices $\rightarrow$ Create from Visit
    \item Select a completed visit --- consultation fee and lab items are added automatically
\end{enumerate}

\section{Recording Payments}

Open an invoice $\rightarrow$ Record Payment. Supported methods: cash, card, online, insurance.

\section{Patient Billing View}

Each patient profile shows their full billing history, totals, and outstanding balance.

\textbf{Patients} can view their own bills via \textbf{My Bills} in the patient portal.

%==============================================================================
\chapter{Pharmacy}
%==============================================================================

\textbf{Who:} Pharmacist, Admin

\begin{itemize}[noitemsep]
    \item \textbf{Inventory} --- drugs, batches, expiry dates
    \item \textbf{Dispense} --- fulfill prescriptions
    \item \textbf{Purchase Orders} --- restock from suppliers
\end{itemize}

Drugs can be bulk-imported via \textbf{Import Data} $\rightarrow$ Pharmacy Drugs.

%==============================================================================
\chapter{Laboratory}
%==============================================================================

\textbf{Who:} Lab Technician, Admin

\begin{itemize}[noitemsep]
    \item \textbf{Test Requests} --- orders from clinical visits
    \item \textbf{Results} --- enter and publish results
\end{itemize}

Lab tests can be imported via CSV under \textbf{Import Data}.

%==============================================================================
\chapter{Human Resources}
%==============================================================================

\textbf{Who:} Admin (attendance, payroll); All staff (leave)

\begin{itemize}[noitemsep]
    \item \textbf{Attendance} --- record daily check-in/out
    \item \textbf{Leave Requests} --- staff submit; admin approves/rejects
    \item \textbf{Payroll} --- monthly payroll runs
\end{itemize}

%==============================================================================
\chapter{Staff Management}
%==============================================================================

\textbf{Who:} Hospital Admin only

\begin{enumerate}[noitemsep]
    \item Sidebar $\rightarrow$ \textbf{Staff}
    \item Click \textbf{Add Staff} --- set username, role, email, password
    \item Available roles: Doctor, Nurse, Receptionist, Accountant, Pharmacist, Lab Technician
\end{enumerate}

\begin{warningbox}
Staff you create belong to \textbf{your hospital only}. They cannot access other hospitals' data. Each staff member signs in with your hospital code at login.
\end{warningbox}

Staff can also be bulk-imported via \textbf{Import Data} $\rightarrow$ Staff (CSV).

%==============================================================================
\chapter{Reports}
%==============================================================================

\textbf{Who:} Admin, Accountant

Available reports include:
\begin{itemize}[noitemsep]
    \item Daily OPD summary
    \item Revenue analysis
    \item Stock levels
\end{itemize}

%==============================================================================
\chapter{Hospital Settings}
%==============================================================================

\textbf{Who:} Admin

Configure:
\begin{itemize}[noitemsep]
    \item Hospital name, address, phone
    \item Brand colors (primary, accent)
    \item Tax rate
    \item Receipt header and footer text
\end{itemize}

%==============================================================================
\chapter{Role Quick Reference}
%==============================================================================

\begin{longtable}{@{}p{3.5cm}p{9cm}@{}}
\toprule
\textbf{Role} & \textbf{Main Menu Items} \\
\midrule
\endhead
Admin &
Dashboard, Import, Patients, Appointments, Clinical, Lab, Pharmacy, Billing, HR, Reports, Staff, Settings \\
Doctor &
Dashboard, My Appointments, My Availability, Patients, Visits, Leave \\
Nurse &
Dashboard, Patients, Wards, Vitals, Leave \\
Receptionist &
Dashboard, Import, Queue, Appointments, Register Patient, Leave \\
Accountant &
Dashboard, Invoices, Payments, Reports, Leave \\
Pharmacist &
Dashboard, Import, Inventory, Dispense, Purchase Orders, Leave \\
Lab Technician &
Dashboard, Import, Test Requests, Results, Leave \\
Patient &
Dashboard, My Appointments, My Records, My Bills \\
\bottomrule
\end{longtable}

%==============================================================================
\chapter{Import Data Reference}
%==============================================================================

\section{Patients CSV (required columns)}

\begin{table}[h]
\centering
\small
\begin{tabular}{@{}lll@{}}
\toprule
\textbf{Column} & \textbf{Required} & \textbf{Example} \\
\midrule
cnic & Yes & 35201-1234567-1 \\
first\_name & Yes & Ahmed \\
last\_name & Yes & Hassan \\
phone & Yes & 03001234567 \\
date\_of\_birth & No & 1990-05-15 \\
gender & No & M / F / O \\
city & No & Islamabad \\
\bottomrule
\end{tabular}
\end{table}

\section{Staff CSV (required columns)}

\texttt{username}, \texttt{first\_name}, \texttt{last\_name}, \texttt{email}, \texttt{role}

Roles: \texttt{doctor}, \texttt{nurse}, \texttt{receptionist}, \texttt{accountant}, \texttt{pharmacist}, \texttt{lab\_tech}

%==============================================================================
\chapter{Troubleshooting}
%==============================================================================

\section{I see another hospital's data}

Sign out. Sign in again with \textbf{your hospital code}. Never share the same \texttt{admin} login across hospitals without specifying the code.

\section{No appointment time slots}

\begin{enumerate}[noitemsep]
    \item Confirm the doctor is \textbf{On Duty}
    \item Confirm weekly schedule is configured
    \item Check the selected date is a working day
\end{enumerate}

\section{Import failed}

\begin{itemize}[noitemsep]
    \item Use UTF-8 CSV (Excel: Save As $\rightarrow$ CSV UTF-8)
    \item Match column headers to the sample file
    \item CNIC must be unique per hospital
    \item Check error table for row numbers
\end{itemize}

\section{Module not available}

Your subscription plan may not include that module. Contact platform support to upgrade.

\section{Account suspended}

Your trial or subscription expired. Contact your administrator or platform support.

%==============================================================================
\chapter{Getting Help}
%==============================================================================

For technical issues, contact your hospital administrator or CARE-U platform support with:
\begin{itemize}[noitemsep]
    \item Your hospital code
    \item Your role and username
    \item Steps to reproduce the issue
    \item Screenshot if applicable
\end{itemize}

\end{document}
